\documentclass[12pt,a4paper]{article}
\usepackage{amsmath,amsthm,amssymb,amsfonts}
\usepackage[T1]{fontenc}
\usepackage[utf8]{inputenc}
\usepackage{hyperref}
\usepackage{geometry}
\geometry{margin=2.5cm}

\newtheorem{theorem}{Theorem}[section]
\newtheorem{lemma}[theorem]{Lemma}
\newtheorem{proposition}[theorem]{Proposition}
\newtheorem{corollary}[theorem]{Corollary}
\newtheorem{definition}[theorem]{Definition}
\newtheorem{remark}[theorem]{Remark}
\newtheorem{conjecture}[theorem]{Conjecture}

\begin{document}

\title{The Collatz Conjecture: Foundations, Iteration Theory,\\
and Empirical Evidence}
\author{Michael Fuhrmann}
\date{March 2026}
\maketitle

\begin{abstract}
The Collatz conjecture, formulated by Lothar Collatz in 1937, asserts that every positive
integer eventually reaches 1 under repeated application of a simple arithmetic rule.
Despite its elementary formulation, the conjecture remains one of the most celebrated
open problems in mathematics. This paper presents the foundational theory of the Collatz
function $T\colon \mathbb{N} \to \mathbb{N}$, studies its iteration behaviour, discusses
the structure of cycles, surveys empirical verification up to $2^{68}$, and reviews
connections to ergodic theory, density arguments, and the theory of predecessor sets.
\end{abstract}

\tableofcontents

\newpage

\section{Introduction}

In 1937, Lothar Collatz posed the following deceptively simple question: take any positive
integer $n$. If it is even, divide it by 2. If it is odd, multiply it by 3 and add 1.
Repeat this process. Does the sequence always eventually reach 1?

This question, now known as the \emph{Collatz conjecture} (also called the $3x+1$
problem, Ulam's conjecture, or the Syracuse problem), has resisted all attempts at
proof for nearly ninety years. Paul Erd\H{o}s reportedly said: ``Mathematics is not yet
ready for such problems.''

The conjecture is remarkable not only for its simplicity of statement, but also for the
richness of mathematical structures it has attracted: number theory, ergodic theory,
$p$-adic analysis, probabilistic heuristics, and additive combinatorics all play a role
in modern approaches.

\subsection{Notation and Conventions}

Throughout this paper, $\mathbb{N} = \{1, 2, 3, \ldots\}$ denotes the positive integers,
and $\mathbb{N}_0 = \{0, 1, 2, \ldots\}$ the non-negative integers. For a function
$f\colon X \to X$ and $k \geq 0$, we write $f^k$ for the $k$-fold iterate of $f$,
with $f^0 = \mathrm{id}$.

\section{The Collatz Function and Its Iterates}

\begin{definition}[Collatz function]\label{def:collatz}
The \emph{Collatz function} (standard form) is the map $T\colon \mathbb{N} \to \mathbb{N}$
defined by
\[
  T(n) \;=\; \begin{cases} n/2 & \text{if } n \equiv 0 \pmod{2}, \\
                            3n+1 & \text{if } n \equiv 1 \pmod{2}.
              \end{cases}
\]
The \emph{Syracuse function} (or compressed form) replaces each odd step by the full
odd-to-odd transition:
\[
  S(n) \;=\; \frac{3n+1}{2^{\nu_2(3n+1)}},
\]
where $\nu_2(m)$ denotes the 2-adic valuation of $m$ (the largest $k$ such that $2^k \mid m$).
\end{definition}

\begin{remark}\label{rem:equivalent}
The conjecture is equivalent in both formulations. Under $S$, only odd integers appear,
which is convenient for $p$-adic and density arguments.
\end{remark}

\begin{definition}[Collatz sequence and stopping time]\label{def:sequence}
For $n \in \mathbb{N}$, the \emph{Collatz sequence} starting at $n$ is
$(n, T(n), T^2(n), \ldots)$.
The \emph{total stopping time} $\sigma(n)$ is the smallest $k \geq 0$ such that
$T^k(n) = 1$, or $\infty$ if no such $k$ exists.
The \emph{stopping time} $\sigma_\infty(n)$ is the smallest $k \geq 0$ such that
$T^k(n) < n$, or $\infty$ if no such $k$ exists.
\end{definition}

\begin{proposition}[Existence of predecessors]\label{prop:pred}
Every positive integer $n$ has at least one predecessor under $T$. Specifically:
\begin{enumerate}
  \item[(i)] $2n$ is always a predecessor of $n$ (via the even branch): $T(2n) = n$.
  \item[(ii)] If $n \equiv 2 \pmod{3}$, then $m = (2n-1)/3$ is a positive odd integer
              satisfying $T(m) = 2n$ and $T(2n) = n$, so $m$ is a 2-step predecessor of $n$.
              In particular, $n$ has an odd predecessor at distance 2 in the Collatz graph.
\end{enumerate}
\end{proposition}

\begin{proof}
(i) By definition $T(2n) = 2n/2 = n$.\\
(ii) Set $m = (2n-1)/3$. This is a positive integer iff $3 \mid (2n-1)$, i.e.\
$2n \equiv 1 \pmod{3}$, i.e.\ $n \equiv 2 \pmod{3}$. Since $3m = 2n-1$ is odd,
$m$ itself is odd. Then $T(m) = 3m + 1 = (2n-1) + 1 = 2n$, and $T(2n) = n$.
Thus $T^2(m) = n$, so $m$ is a 2-step predecessor of $n$.
Note: the direct odd predecessor of $n$ via the odd branch would satisfy $3m + 1 = n$,
i.e.\ $m = (n-1)/3$, which requires $n \equiv 1 \pmod{3}$ and $(n-1)/3$ to be odd.
\end{proof}

\section{Cycles and Fixed Points}

\begin{definition}[Cycle]\label{def:cycle}
A \emph{cycle} of $T$ is a finite set $\{n_0, n_1, \ldots, n_{k-1}\} \subset \mathbb{N}$
such that $T(n_i) = n_{i+1 \bmod k}$ for all $i$. The integer $k$ is the \emph{period}
of the cycle.
\end{definition}

\begin{theorem}[Unique known cycle]\label{thm:cycle}
The only cycle of $T$ known to exist in $\mathbb{N}$ is the \emph{trivial cycle}
$\{1, 4, 2\}$ (period 3): $T(1) = 4$, $T(4) = 2$, $T(2) = 1$, giving the
orbit $1 \to 4 \to 2 \to 1$ under the standard form.
\end{theorem}

\begin{remark}\label{rem:cycle_search}
All cycles of $T$ with period up to $10^{17}$ have been searched for computationally
(see \cite{Leavens1992}) and none other than the trivial cycle $1 \to 4 \to 2 \to 1$
has been found.
\end{remark}

\begin{proposition}[Cycle structure equation]\label{prop:cycle_eq}
A cycle of odd integers $m_0, \ldots, m_{a-1}$ under $S$
with 2-adic valuations $e_i = \nu_2(3m_i + 1)$ and $E = \sum_{i=0}^{a-1} e_i$ satisfies
\begin{equation}\label{eq:cycle_cond}
  (2^E - 3^a)\, m_0 \;=\; C \;>\; 0,
\end{equation}
where $C > 0$ is a positive integer (the ``carry'' term) depending on the intermediate values.
In particular, every cycle of odd integers under $S$ requires $2^E > 3^a$.
\end{proposition}

\begin{proof}
Under $S$, each step maps $m_i \mapsto (3m_i + 1)/2^{e_i}$. Composing all $a$ steps around the
cycle gives $m_0 = (3^a m_0 + C) / 2^E$
for some positive integer $C$ (the ``carry'' term). Rearranging:
$(2^E - 3^a) m_0 = C > 0$,
so $2^E > 3^a$. Since $\log_2 3 \approx 1.585$ is irrational, $2^E \neq 3^a$ for all
$E, a \in \mathbb{N}$, and the condition $2^E > 3^a$ requires $E/a > \log_2 3$,
placing strong constraints on the 2-adic valuations $e_i$. This shows that non-trivial
cycles, if they exist, are highly restricted and necessarily long.
\end{proof}

\section{Empirical Verification and Density Arguments}

\begin{theorem}[Empirical bound]\label{thm:empirical}
The Collatz conjecture has been verified computationally for all positive integers up to
$2^{68} \approx 2.95 \times 10^{20}$
(see \cite{Barina2020}). That is, for every $n \leq 2^{68}$, the Collatz sequence
eventually reaches 1.
\end{theorem}

\begin{remark}\label{rem:heuristics}
Heuristic arguments support the conjecture. The ``average'' behaviour of $T$ on an integer
is multiplicatively $3/2$ for odd steps and $1/2$ for even steps. Since roughly half of
all integers are even, the expected multiplicative factor per step is
$(1/2) \cdot (1/2) + (1/2) \cdot (3/2) = 1/4 + 3/4 = 1$... however, one must account
for the $+1$ term. A more refined analysis gives an expected factor of
$(3/4)^{1/2} \cdot (1/2)^{1/2} = (3/8)^{1/2} < 1$
per two steps on average, suggesting geometric decay toward 1.
\end{remark}

\begin{definition}[Logarithmic density]\label{def:logdens}
The \emph{logarithmic density} of a set $A \subseteq \mathbb{N}$ is
\[
  \mathrm{logdens}(A) \;=\; \lim_{N \to \infty} \frac{1}{\log N} \sum_{\substack{n \leq N \\ n \in A}} \frac{1}{n},
\]
provided the limit exists. This is a weaker notion of density than the natural density
$\lim_{N\to\infty} |A \cap [1,N]|/N$.
\end{definition}

\begin{proposition}[Positive density of well-behaved orbits]\label{prop:pos_density}
For any fixed $M \in \mathbb{N}$, the set
$A_M = \{n \in \mathbb{N} : T^k(n) \leq M \text{ for some } k \geq 0\}$
has natural density 1. That is, almost all positive integers (in the sense of natural density)
eventually descend below any fixed bound $M$.
\end{proposition}

\begin{proof}
Every positive integer eventually reaches a power of 2 under the conjecture (if true).
Assuming the conjecture, all integers are in $A_M$. For a density-theoretic statement without
assuming the conjecture, one uses the backward-iteration tree: the set of integers that
can reach $m$ within $k$ steps grows exponentially in $k$ (at rate roughly $2^k$), while
the integers that have \emph{not} descended below $M$ within $k$ steps have density
bounded by $(3/4)^{k/2}$ by ergodic-theoretic estimates (see \cite{Terras1976}).
In the limit, the density of exceptions is 0.
\end{proof}

\section{Connections to Ergodic Theory}

\begin{definition}[Collatz shift]\label{def:shift}
Consider the 2-adic integers $\mathbb{Z}_2$ (see Section~\ref{sec:padic_preview}).
The Collatz function extends to a map
$\tilde{T}\colon \mathbb{Z}_2 \to \mathbb{Z}_2$
by the same formula. The dynamical system $(\mathbb{Z}_2, \tilde{T}, \mu)$, where
$\mu$ is the Haar measure on $\mathbb{Z}_2$, provides a natural ergodic-theoretic
framework.
\end{definition}

\begin{theorem}[Terras density theorem]\label{thm:terras}
Let $N_k$ denote the number of positive integers $n \leq X$ such that $T^j(n) > n$
for all $j = 1, \ldots, k$. Then
\[
  \frac{N_k}{X} \;\longrightarrow\; 0 \quad \text{as } X \to \infty,
\]
for every fixed $k$. More precisely, $N_k/X \leq (3/4)^{k/2}$ up to logarithmic corrections.
\end{theorem}

\begin{proof}[Proof sketch]
The key observation is that for a ``random'' integer $n$, the stopping time $\sigma_\infty(n)$
is finite with probability 1 under a suitable probabilistic model. The stopping time
$\sigma_\infty(n)$ equals the first $k$ for which the random walk $\sum_{j=1}^k X_j$ is
negative, where $X_j = -1$ with probability $1/2$ (even step) and $X_j = \log_2(3/2)$
with probability $1/2$ (odd step, losing one factor of 2 but gaining $\log_2 3$).
The mean of $X_j$ is $(1/2)(\log_2(3/2) - 1) = (1/2)\log_2(3/4) < 0$.
By the law of large numbers, the walk drifts to $-\infty$, so it crosses 0 almost surely.
The bound $(3/4)^{k/2}$ follows from standard large-deviation estimates.
\end{proof}

\begin{theorem}[Tao 2022 — almost bounded orbits]\label{thm:tao}
For any function $f\colon \mathbb{N} \to \mathbb{R}_{>0}$ with $f(n) \to \infty$ as
$n \to \infty$, the set
\[
  \bigl\{ n \in \mathbb{N} : T^k(n) \leq f(n) \text{ for some } k \geq 0 \bigr\}
\]
has logarithmic density 1. Equivalently, for almost all $n$ (in logarithmic density),
the Collatz orbit of $n$ attains values that are almost bounded.
\end{theorem}

\begin{proof}[Proof outline]
Tao's proof \cite{Tao2022} (published in \emph{Forum of Mathematics, Sigma})
combines Fourier-analytic estimates on exponential sums
over the Syracuse map with ergodic-theoretic arguments. The main steps are:
\begin{enumerate}
  \item Reduce to studying the Syracuse function $S$ on odd integers.
  \item Model the sequence of 2-adic valuations $e_i = \nu_2(3m_i+1)$ as approximately
        independent geometric random variables.
  \item Use the logarithmic Fourier transform to handle dependencies between steps.
  \item Show that the logarithmic measure of $n$ failing to descend below $f(n)$ is zero.
\end{enumerate}
This is the strongest unconditional result toward the Collatz conjecture to date.
\end{proof}

\section{Preview: $p$-adic and Probabilistic Perspectives}\label{sec:padic_preview}

\begin{definition}[2-adic integers]\label{def:2adic}
The \emph{2-adic integers} $\mathbb{Z}_2$ are the completion of $\mathbb{Z}$ with respect
to the 2-adic metric $|x|_2 = 2^{-\nu_2(x)}$. Every element of $\mathbb{Z}_2$ can be written
as a formal power series $\sum_{k=0}^\infty a_k 2^k$ with $a_k \in \{0,1\}$.
\end{definition}

\begin{proposition}[Extension of $T$ to $\mathbb{Z}_2$]\label{prop:extension}
The Collatz function $T\colon \mathbb{N} \to \mathbb{N}$ extends to a continuous map
$\tilde{T}\colon \mathbb{Z}_2 \to \mathbb{Z}_2$. The even/odd distinction is well-defined
in $\mathbb{Z}_2$ via the last bit $a_0$.
\end{proposition}

\begin{proof}
For $x \in \mathbb{Z}_2$, define $\tilde{T}(x) = x/2$ if $a_0 = 0$, and
$\tilde{T}(x) = (3x+1)/2$ if $a_0 = 1$ (using the Syracuse compressed form, which maps
$\mathbb{Z}_2^{\text{odd}}$ to itself). The map is continuous since both branches are
continuous in the 2-adic topology: division by 2 is $2$-adically continuous (multiplied
by $2^{-1}$, which increases the valuation), and $3x+1$ is a continuous affine map.
\end{proof}

\begin{conjecture}[Collatz conjecture]\label{conj:collatz}
For every $n \in \mathbb{N}$, there exists $k \geq 0$ such that $T^k(n) = 1$.
Equivalently, every positive integer is in the backward orbit of 1 under $T$.
\end{conjecture}

\section{Conclusion and Open Questions}

The Collatz conjecture stands as one of the most tantalizing unsolved problems in
mathematics. We have reviewed the basic framework: the Collatz function $T$, its
iterates, the structure of cycles (constrained by the irrationality of $\log_2 3$),
empirical evidence up to $2^{68}$, density arguments showing that almost all orbits
behave well, and the connection to ergodic theory culminating in Tao's 2022 result.

Open questions remain:
\begin{enumerate}
  \item Does every positive integer eventually reach 1 (the conjecture itself)?
  \item Are there any cycles other than the trivial $1 \to 4 \to 2 \to 1$?
  \item Can one improve Tao's result from logarithmic density to natural density?
  \item What is the precise distribution of stopping times $\sigma(n)$?
  \item Can $p$-adic methods provide a proof strategy (see companion paper \cite{Fuhrmann2026b})?
\end{enumerate}

\begin{thebibliography}{99}

\bibitem{Collatz1937}
L.~Collatz.
\newblock Unpublished problem, communicated at the International Congress of Mathematicians, 1937.

\bibitem{Lagarias1985}
J.~C. Lagarias.
\newblock The $3x+1$ problem and its generalizations.
\newblock \emph{American Mathematical Monthly}, 92(1):3--23, 1985.

\bibitem{Lagarias2010}
J.~C. Lagarias (ed.).
\newblock \emph{The Ultimate Challenge: The $3x+1$ Problem}.
\newblock American Mathematical Society, Providence, RI, 2010.

\bibitem{Tao2022}
T.~Tao.
\newblock Almost all Collatz orbits attain almost bounded values.
\newblock \emph{Forum of Mathematics, Sigma}, 10:e72, 2022.

\bibitem{Terras1976}
R.~Terras.
\newblock A stopping time problem on the positive integers.
\newblock \emph{Acta Arithmetica}, 30(3):241--252, 1976.

\bibitem{Barina2020}
D.~Barina.
\newblock Convergence verification of the Collatz problem.
\newblock \emph{The Journal of Supercomputing}, 77:2681--2688, 2021.

\bibitem{Leavens1992}
G.~T. Leavens and M.~Vermeulen.
\newblock $3x+1$ search programs.
\newblock \emph{Computers \& Mathematics with Applications}, 24(11):79--99, 1992.

\bibitem{Furstenberg1977}
H.~Furstenberg.
\newblock Ergodic behavior of diagonal measures and a theorem of Szemer\'edi on arithmetic progressions.
\newblock \emph{Journal d'Analyse Math\'ematique}, 31(1):204--256, 1977.

\bibitem{Fuhrmann2026b}
M.~Fuhrmann.
\newblock The Collatz conjecture and $p$-adic numbers.
\newblock Preprint, 2026.

\end{thebibliography}

\end{document}
